\chapter{Introduction}
\label{sec:intro}

Blockchain technology provides decentralized, robust, and programmable ledgers
at Internet scale. Developers can deploy smart contracts onto a blockchain
to encode arbitrarily complicated transaction rules that operate on the ledger. 
%The faithful execution of smart contracts are enforced and verified by all
%participants of the blockchain network. 
Blockchains and smart contracts have become the key infrastructure for various
decentralized financial services (DeFi). The Total Value Locked (TVL) in $1417$
DeFi smart contracts has reached $85.64$ billion USD on June 1st
2022~\cite{DeFillama}.


Security attacks are critical threats to smart contracts. Attackers can send
malicious transactions to exploit vulnerabilities of a contract to steal
millions of dollars from the contract users. Particularly, a new kind of
security attacks that use large amount of digital assets to exploit design
flaws in DeFi contracts has become prevalent. Because such attacks typically
borrow the used assets from flash loan contracts, they are called \emph{flash
loan attacks}~\cite{SlowMistStats, mckay2022defi}. Among the top $200$ costliest
attacks recorded in Rekt Database,  the financial loss caused by $36$
flash loan attacks exceeds $418$ million USD~\cite{mckay2022defi}.

%The flash loan is a special kind of loan that occurs inside one transaction,
%i.e., the borrower must return all borrowed assets back before the end of the
%transaction otherwise the flash loan protocol will abort the transaction to
%cancel the loan.
A malicious flash loan attack transaction typically contains a sequence of
actions (i.e., function calls to smart contracts). The first action borrows a
very large sum of digital assets from a flash loan contract and the last action
returns the borrowed assets. The sequence of actions in the middle interacts
with multiple DeFi contracts using the borrowed assets to exploit their design
flaws. When a DeFi contract fails to consider corner cases caused by the large
sum of borrowed assets, the attacker may extract prohibitive profits. For
example, many flash loan attacks use borrowed assets to temporarily manipulate
asset prices in a DeFi contract to fool the contract to make unfavorable trades
with the attacker~\cite{qin2021attacking,cao2021flashot}. Although researchers
have developed many automated program analysis and verification
techniques~\cite{mossberg2019manticore,feng2019precise, grieco2020echidna,
brent2020ethainter} to detect and eliminate bugs in smart
contracts, these techniques cannot handle flash loan attack vulnerabilities.
This is because such vulnerabilities are design flaws rather than
implementation bugs. Moreover, these techniques typically operate with one
contract at a time, but flash loan attacks almost always involve multiple DeFi
contracts interacting with each other.


\noindent \textbf{\name:}
We present {\name}, the first automated end-to-end program synthesis
tool for detecting flash loan attack vulnerabilities. Given a set of smart
contracts and candidate actions in these contracts, {\name} automatically
synthesizes an action sequence along with all action parameters to interact
with the contracts to exploit potential flash loan vulnerabilities.

The main challenge {\name} faces is that the underlying logic of DeFi actions
is often too sophisticated for standard solvers or optimizers to handle. Even
if the action sequence was already known, a naive application of symbolic
execution might not be able to find action parameters because it may need to extract
overly complicated symbolic constraints causing the solvers to
time out. Moreover, {\name} synthesizes the action sequence and the action
parameters together and therefore faces an additional search space explosion
challenge.

{\name} addresses these challenges with its novel
\emph{synthesis-via-approximation} technique. Instead of attempting to extract
accurate symbolic expressions from smart contract code, {\name} collects data
points to approximate the effect of contract actions with polynomials and
interpolations. {\name} then uses the approximated expressions to drive the
synthesis. {\name} also incrementally improves the approximation with its novel
\emph{counterexample-driven approximation refinement} techniques, i.e., if the
synthesis fails because of a large deviation caused by the approximations,
{\name} collects the corresponding data points as counterexamples to
iteratively refine the approximations. The combination of these techniques
allows the underlying optimizer of {\name} to work with more tractable
expressions. It also decouples the two difficult tasks, finding the action
sequence and finding the action parameters. When working with a set of
coarse-grained approximated expressions, {\name} can filter out unproductive
action sequences with a small cost.

\noindent \textbf{Experimental Results:} 
We evaluate {\name} on $16$ DeFi benchmark protocols that were victims to flash
loan attacks and $2$ DeFi benchmark protocols from Damn Vulnerable DeFi
challenges~\cite{DVD}. {\name} automatically synthesizes adversarial attacks for $16$ out
of the 18 benchmarks. For comparison, a baseline that operates manually crafted
accurate action summaries only synthesizes attacks for 7 out of the 18.

\noindent \textbf{Contributions:}
This thesis makes the following contributions:
\begin{itemize}
\item \textbf{\name:} The first automated
end-to-end program synthesis tool for detecting flash loan attack
vulnerabilities. It enables approximate attack synthesis 
without diving into sophisticated logics of DeFi contracts. 

\item \textbf{Synthesis via Approximation:} A novel synthesis-via-approximation technique to handle
        sophisticated logics of DeFi contracts. 

\item \textbf{Counterexample Driven Approximation Refinement:} A novel counterexample driven approximation refinement technique to
        incrementally improve the approximation during the synthesis process.
\item \textbf{Experimental Evaluation:} 
We implemented {\name} in a tool and evaluated it on $16$ protocols that were victims to flash
loan attacks and $2$ fictional flash loan attacks. 
\end{itemize}

The rest of the thesis is organized as follows. Section~\ref{sec:background} briefly introduces background information. 
Section~\ref{sec:example} 
presents a motivating example to give a high-level overview of {\name}. 
In Section ~\ref{sec:preli}, we define the problem we tackle in this thesis. 
Section~\ref{sec:approach} presents our core synthesis algorithm.
%and 
%the counterexample driven approximation refinement technique. 
Section~\ref{sec:implementation} describes {\name}, an implementation of our techniques. 
Section~\ref{sec:evaluation} evaluates {\name} on $18$ benchmarks. 
In Sections~\ref{sec:related} and \ref{sec:threats}, 
we discuss related work and threats to validity, respectively.  
Section~\ref{sec:conclusion} concludes the thesis.
